\chapter{Separability and Entanglement}
\label{ch:entanglement}

Chapter~\ref{ch:partial-trace} left one question open deliberately: is
$\rho(c)$, for $c\ne0$, entangled, or merely classically correlated in some
other way invisible to marginals? This chapter answers it — and shows the
answer depends on which scalar field is doing the composing, which is exactly
the kind of question this book is positioned to ask.

\section{Separability Relative to a Scalar Field}

\begin{definition}[Separability]
\label{def:separability}
\index{Separability (definition)}
Fix $\mathbb F \in \{\R,\C\}$. A state $\rho \in \mathcal D_{d_Ad_B}(\mathbb
F)$ is \emph{$\mathbb F$-separable} if it is a convex combination of
$\mathbb F$-product states,
\begin{equation}
\rho = \sum_k p_k\, \rho_A^{(k)} \otimes \rho_B^{(k)}, \qquad p_k \ge 0,\
\sum_k p_k = 1,\qquad \rho_A^{(k)} \in \mathcal D_{d_A}(\mathbb F),\ \rho_B^{(k)}
\in \mathcal D_{d_B}(\mathbb F).
\end{equation}
A state that is not $\mathbb F$-separable is \emph{$\mathbb F$-entangled}.
\end{definition}

The qualifier is essential: $\mathcal D_d(\R) \subsetneq \mathcal D_d(\C)$
(Chapters~\ref{ch:local-tomography}--\ref{ch:positivity}), so the set of
available local product states differs between the two theories, and
separability is a property of a state \emph{relative to} which local states
are admitted — not an intrinsic property of a matrix alone.

\section{The \texorpdfstring{$J\otimes J$}{J⊗J} Witness Over \texorpdfstring{$\R$}{R}}

\begin{lemma}
\label{lem:trace-symmetric-antisymmetric}
For any real symmetric $\rho_A \in \mathrm{Herm}_\R(d)$, $\operatorname{tr}(\rho_A
J) = 0$ whenever $J^\top = -J$ is antisymmetric (same dimension).
\end{lemma}

\begin{proof}
$\operatorname{tr}(\rho_A J) = \operatorname{tr}\big((\rho_A J)^\top\big) =
\operatorname{tr}(J^\top \rho_A^\top) = \operatorname{tr}(-J\rho_A) =
-\operatorname{tr}(J\rho_A) = -\operatorname{tr}(\rho_A J)$ (using
cyclicity of trace), so $\operatorname{tr}(\rho_A J) = 0$.
\end{proof}

\begin{theorem}[Real entanglement witness]
\label{thm:real-witness}
\index{Real entanglement witness (theorem)}
For every $\R$-separable two-rebit state $\rho$, $\operatorname{tr}\big[\rho\,
(J\otimes J)\big] = 0$.
\end{theorem}

\begin{proof}
For a product state, $\operatorname{tr}\big[(\rho_A\otimes\rho_B)(J\otimes
J)\big] = \operatorname{tr}(\rho_A J)\,\operatorname{tr}(\rho_B J) = 0\cdot 0 =
0$ by Lemma~\ref{lem:trace-symmetric-antisymmetric} (using the standard
identity $\operatorname{tr}[(A\otimes B)(C\otimes D)] = \operatorname{tr}(AC)
\operatorname{tr}(BD)$). By linearity of trace, the same holds for any convex
combination.
\end{proof}

\begin{proposition}
\label{prop:c-classification}
$\operatorname{tr}\big[\rho(c)(J\otimes J)\big] = c$.
\end{proposition}

\begin{proof}
$\operatorname{tr}\big[\rho(c)(J\otimes J)\big] = \tfrac14\big(
\operatorname{tr}(J\otimes J) + c\,\operatorname{tr}\big[(J\otimes J)^2\big]
\big) = \tfrac14\big(0 + c\cdot 4\big) = c$, using
$\operatorname{tr}(J\otimes J)=0$ and $(J\otimes J)^2 = I_4$ (Chapter
\ref{ch:tensor-product}), so $\operatorname{tr}[(J\otimes J)^2] =
\operatorname{tr}(I_4) = 4$.
\end{proof}

\begin{corollary}
\label{cor:real-classification}
$\rho(c)$ is $\R$-separable if and only if $c=0$.
\end{corollary}

\begin{proof}
If $c\ne0$, Theorem~\ref{thm:real-witness} and Proposition
\ref{prop:c-classification} together rule out $\R$-separability. If $c=0$,
$\rho(0) = \tfrac12I \otimes \tfrac12I$ is manifestly a real product state
(indeed $\tfrac12I \in \mathcal D_2(\R)$), hence separable.
\end{proof}

This is the four-chapter chain complete: $J\otimes J$ appears as an extra
basis vector invisible to local product operators (Chapter
\ref{ch:tensor-product}), becomes a bounded physical coordinate $|c|\le1$
under positivity (Chapter~\ref{ch:positivity}), vanishes under both
marginals (Chapter~\ref{ch:partial-trace}), and is now shown to be exactly an
entanglement witness: the same coefficient invisible to every local
measurement is precisely what makes $\rho(c)$, $c\ne0$, entangled.

\section{Complexification: The Same States, a Different Theory}

\begin{proposition}
\label{prop:JJ-is-sigmay}
$J = -i\sigma_y$, so $J\otimes J = -\sigma_y\otimes\sigma_y$.
\end{proposition}

This is immediate from $\sigma_y = \begin{psmallmatrix}0&-i\\i&0
\end{psmallmatrix}$. The point is that $J\otimes J$, unavailable as a product
of \emph{real} local observables (Chapter~\ref{ch:tensor-product}), is
exactly $-1$ times a product of perfectly ordinary \emph{complex} local
observables, $\sigma_y\otimes\sigma_y$.

\begin{theorem}[The same family is entirely $\C$-separable]
\label{thm:complex-separable}
\index{The same family is entirely C-separable (theorem)}
$\rho(c)$ is $\C$-separable for every $c\in[-1,1]$.
\end{theorem}

\begin{proof}
Let $|{\pm}y\rangle$ be the $\pm1$-eigenvectors of $\sigma_y$. Since $J\otimes
J = -\sigma_y\otimes\sigma_y$, and $\sigma_y\otimes\sigma_y$ has eigenvalue
$+1$ on $\{|{+}y,{+}y\rangle,|{-}y,{-}y\rangle\}$ and eigenvalue $-1$ on
$\{|{+}y,{-}y\rangle,|{-}y,{+}y\rangle\}$ (products of $\sigma_y$-eigenvalues),
direct substitution gives the exact identity
\begin{equation}
\rho(c) = \frac{1+c}{4}\,|{+}y,{-}y\rangle\langle{+}y,{-}y| +
\frac{1+c}{4}\,|{-}y,{+}y\rangle\langle{-}y,{+}y| +
\frac{1-c}{4}\,|{+}y,{+}y\rangle\langle{+}y,{+}y| +
\frac{1-c}{4}\,|{-}y,{-}y\rangle\langle{-}y,{-}y|.
\end{equation}
Each term is a pure complex product state; the four weights are nonnegative
and sum to $1$ exactly when $|c|\le1$, which always holds
(Chapter~\ref{ch:positivity}). This exhibits $\rho(c)$ as an explicit convex
combination of $\C$-product states for every admissible $c$.
\end{proof}

\begin{remark}
This is stronger than a partial or asymptotic statement: the \emph{entire}
one-parameter family that is $\R$-entangled for every $c\ne0$
(Corollary~\ref{cor:real-classification}) is $\C$-separable across its whole
range. Nothing about $\rho(c)$ itself changes — it is the same real $4\times4$
matrix throughout — only the theory doing the composing. Separability is not
an intrinsic property of a state; it is a property of a state \emph{relative
to} the field over which local product states are admitted.
\end{remark}

\begin{scopebox}
\scopetitle{}This is a precise claim about one family of matrices under
Definition~\ref{def:separability}, not a general verdict that "entanglement
isn't physically real." It does not say every entangled state in ordinary
complex quantum mechanics becomes separable under some other field, and it
does not say separability judgments are arbitrary or a matter of choice
once a scalar field is fixed — within ordinary complex quantum mechanics,
$\rho(c)$'s entanglement status (for whatever states arise there) is
perfectly well-defined. What is field-relative is the comparison
\emph{across} candidate theories with different admissible local product
states, which is the question this book's Parts I--II were built to ask.
\end{scopebox}

\begin{ledgerentry}[End of Chapter~\ref{ch:entanglement}]
\textbf{Established:} $\rho(c)$ is $\R$-entangled for every $c\ne0$, exhibited
by the explicit witness $J\otimes J$ (Corollary~\ref{cor:real-classification});
the identical family is $\C$-separable for every $c\in[-1,1]$, exhibited by an
explicit four-term product-state decomposition
(Theorem~\ref{thm:complex-separable}); separability is therefore
field-relative, not intrinsic to a state.\\
\textbf{Not established:} anything about entanglement in the full generality
of $d>2$ composite systems, multipartite entanglement, entanglement measures
or their operational meaning, or the fate of quaternionic composite states
under this classification (the overcompleteness of Chapter
\ref{ch:tensor-product} suggests the very notion of a $\Hset$-product state
needs care before separability can even be defined there). These remain open.
\end{ledgerentry}
